\documentclass{article}
\usepackage{geometry}
\geometry{a4paper, margin=1in}
\title{ML Research Proposal: Deep Decoding of Spatio-Temporal Bioelectric Patterns}
\author{Candidate Name}
\date{Suggested Supervisors: Prof. Dr. Martin Weigert, Dr. Elias Barriga}

\begin{document}
\maketitle

\section*{Abstract}
The dynamics of morphogenesis are driven by local bioelectric signals ($\text{V}_{\text{mem}}$ gradients) and mechanical forces, both of which are measured via noisy, high-dimensional 3D live-imaging data. Extracting reliable, quantitative features for modeling (such as local cell curvature, velocity fields, and bioelectric pattern boundaries) is a major bottleneck. This project will develop a custom \textbf{Deep Learning Bioimage Analysis Pipeline} to achieve a **Spatio-Temporal Understanding** of the morphogenetic program. Leveraging the expertise in machine learning for spatial understanding \cite{4, 6}, we will train state-of-the-art networks (e.g., U-Nets or Transformer-based models) to: 1) \textbf{Robustly segment and track} individual cells, nuclei, and internal organelles across 4D (3D+time) microscopy data from the Barriga lab (bioelectric/morphogenesis assays). 2) \textbf{Extract quantitative local features}, including the cell-to-cell gradient of membrane voltage ($\text{V}_{\text{mem}}$) and the orientation of the cytoskeleton (the active polarity vector $p_a$). The output will be a set of clean, physics-ready feature maps that can be directly used as input for TopoSPAM models or a downstream predictive classifier, thereby bridging bioimage analysis and biophysics modeling.

\section*{Related Work}
\begin{itemize}
    \item \textbf{Machine Learning for Spatial Understanding:} Advanced bioimage analysis techniques (e.g., CARE, StarDist) for robust segmentation, object detection, and tracking in microscopy data \cite{4, 6}.
    \item \textbf{Bioelectricity in Morphogenesis:} Experimental work demonstrating the role of $\text{V}_{\text{mem}}$ as a key communication signal during tissue shaping (e.g., Levin's work; Barriga Lab data).
    \item \textbf{Active Vertex Models:} Requirement for precise segmentation of cell boundaries and polarity vectors to parameterize discrete cell-level simulations \cite{1}.
\end{itemize}

\end{document}
